\section{Turing Machine --- Build --- Step Block Loop}
%%% generated by the blueprint pipeline from TCSlib.Complexity.TuringMachine.Build.StepBlockLoop
\subsection*{Turing Machine --- Build --- Step Block Loop: Introduction}

Declarations of the Lean module \texttt{TCSlib.Complexity.TuringMachine.Build.StepBlockLoop}.

\subsection*{Turing Machine --- Build --- Step Block Loop: Results}

\begin{definition}[Counter-controlled round machine]
\lean{Turing.StepBlock.loopTM}
\label{Turing.StepBlock.loopTM}
\leanok
\uses{Turing.CounterLoopState, Turing.MultiTapeTM, Turing.StepBlock.Ph, Turing.StepBlock.exit, Turing.StepBlock.roundTM, Turing.counterLoopTM}
For each $k\in\mathbb{N}$, the loop machine is a counter-controlled machine whose body
is the compiled round machine. It uses the designated true and false exit states, with
the false exit supplied as an optional exit, and its control states are the counter-loop
states formed from the compiled machine's state type.
\end{definition}

\begin{definition}[Hosted counter-controlled round machine]
\lean{Turing.StepBlock.hostedLoopTM}
\label{Turing.StepBlock.hostedLoopTM}
\leanok
\uses{Turing.CounterLoopState, Turing.MultiTapeTM, Turing.StepBlock.Ph, Turing.StepBlock.loopTM, Turing.vhostEmitTM}
For each $k\in\mathbb{N}$, the hosted loop machine is obtained by hosting the
counter-controlled round machine and emitting its output. Its work tapes consist of one
additional tape together with the loop machine's tapes, and its states pair the loop
machine's control state with a Boolean value.
\end{definition}

\begin{definition}[First round exit time]
\lean{Turing.StepBlock.loopRoundTime}
\label{Turing.StepBlock.loopRoundTime}
\leanok
\uses{Turing.Cfg, Turing.MultiTapeTM.runFrom, Turing.StepBlock.Ph, Turing.StepBlock.exit, Turing.StepBlock.roundTM}
For a compiled-machine configuration $d$ on a Boolean input word $x$, define its round
time to be the least positive integer $t$ for which the round machine started at $d$ is
in one of the two designated exit states after $t$ steps, if such a time exists; if no
such positive time exists, set the round time to zero.
\end{definition}

\begin{lemma}[Successive loop-orbit configurations]
\lean{Turing.StepBlock.loopOrbit_succ}
\label{Turing.StepBlock.loopOrbit_succ}
\leanok
\uses{Turing.Cfg, Turing.MultiTapeTM.runFrom, Turing.StepBlock.Ph, Turing.StepBlock.loopRoundTime, Turing.StepBlock.roundTM, Turing.counterLoopOrbit, Turing.roundOrbit_succ}
\difficulty{1}
For any compiled-machine configuration $d$ and natural number $r$, the configuration
after $r+1$ exit-delimited rounds equals the result of running the compiled machine from
the configuration after $r$ rounds for exactly the duration assigned to that
configuration's next round.
\end{lemma}

\begin{lemma}[Simulation of One Machine Step]
\lean{Turing.StepBlock.loopRoundTime_spec}
\label{Turing.StepBlock.loopRoundTime_spec}
\leanok
\uses{Turing.Action.apply, Turing.Cfg, Turing.Cfg.inputSymbol, Turing.Cfg.workTapeSymbols, Turing.MultiTapeTM, Turing.MultiTapeTM.runFrom, Turing.MultiTapeTM.step, Turing.StepBlock.Ph, Turing.StepBlock.Rep, Turing.StepBlock.exit, Turing.StepBlock.loopRoundTime, Turing.StepBlock.roundTM, Turing.StepBlock.roundTM_effect, Turing.StepBlock.roundTM_firstExit_unique, Turing.tableBitsK}
\difficulty{5}
Let $k,n\in\mathbb N$, let $x$ be a finite binary input, let $M$ be a $k$-tape machine
with binary work tapes and states indexed by $\{0,\ldots,n\}$, and let $c$ be one of its
configurations whose state is present. Let $d$ be a configuration of the round machine
on the same input, with state at the live entry, and suppose $d$ represents $c$ and the
transition table of $M$. If $\tau$ is the first positive time at which the round machine
reaches either exit, then $0<\tau\le 50(L+1)$, where $L$ is the length of the encoded
transition table. At time $\tau$ the round machine is at the live exit exactly when the
state of the one-step successor of $c$ is present; its configuration represents that
successor and its output is unchanged. At every intermediate positive time before
$\tau$, it is at neither exit.
\end{lemma}

\begin{lemma}[Representation along loop iterations]
\lean{Turing.StepBlock.loopOrbit_rep}
\label{Turing.StepBlock.loopOrbit_rep}
\leanok
\uses{Turing.Cfg, Turing.MultiTapeTM, Turing.MultiTapeTM.runFrom, Turing.MultiTapeTM.runFrom_succ_eq_step', Turing.MultiTapeTM.runFrom_zero, Turing.StepBlock.Ph, Turing.StepBlock.Rep, Turing.StepBlock.exit, Turing.StepBlock.loopOrbit_succ, Turing.StepBlock.loopRoundTime, Turing.StepBlock.loopRoundTime_spec, Turing.StepBlock.roundTM, Turing.counterLoopOrbit}
\difficulty{4}
Let $M$ be a $k$-work-tape binary machine with states indexed by $\text{Fin}(n+1)$, and
let $c$ be an initial configuration with a present state. Suppose a compiled
configuration $d$ starts at the live exit anchor and represents $c$. For a natural
number $r$, assume the machine has a present state at every time $u<r$. Let $e$ be the
configuration obtained by iterating the compiled machine for $r$ rounds, each round
lasting until it reaches an exit anchor. Then $e$ is at the exit anchor determined by
whether the simulated machine is still running after $r$ steps, $e$ represents the
simulated configuration at that time, and $e$ has the same output as $d$.
\end{lemma}

\begin{lemma}[One compiled round reaches its live exit]
\lean{Turing.StepBlock.loopOrbit_round}
\label{Turing.StepBlock.loopOrbit_round}
\leanok
\uses{Turing.Cfg, Turing.MultiTapeTM, Turing.MultiTapeTM.runFrom, Turing.StepBlock.Ph, Turing.StepBlock.Rep, Turing.StepBlock.exit, Turing.StepBlock.loopOrbit_rep, Turing.StepBlock.loopRoundTime, Turing.StepBlock.loopRoundTime_spec, Turing.StepBlock.roundTM, Turing.counterLoopOrbit, Turing.tableBitsK}
\difficulty{3}
Let $M$ be a $k$-work-tape binary machine with states indexed by $\text{Fin}(n+1)$, and
let $c$ and $d$ be configurations such that $d$ starts at the live exit anchor and
represents $c$. For a natural number $r$, suppose the simulated machine has a present
state at every time $u\le r$. Let $e$ be the configuration after
iterating the compiled machine for $r$ exit-delimited rounds. Then $e$ is again at the
live exit anchor, and the duration of the next round is positive and at most
$50(|\text{tableBits}(M)|+1)$. At every strictly intermediate time in that round, the
compiled machine is in a state different from both exit anchors.
\end{lemma}

\begin{theorem}[Bounded Simulation of Repeated Steps]
\lean{Turing.StepBlock.loopTM_run}
\label{Turing.StepBlock.loopTM_run}
\leanok
\uses{Turing.Cfg, Turing.FinTM.bufferTape, Turing.MultiTapeTM, Turing.MultiTapeTM.exists_first_halt, Turing.MultiTapeTM.runFrom, Turing.MultiTapeTM.runFrom_state_ne_none_of_le, Turing.MultiTapeTM.runFrom_zero, Turing.StepBlock.Ph, Turing.StepBlock.Rep, Turing.StepBlock.exit, Turing.StepBlock.loopOrbit_rep, Turing.StepBlock.loopOrbit_round, Turing.StepBlock.loopRoundTime, Turing.StepBlock.loopTM, Turing.StepBlock.roundTM, Turing.counterLoopCfg, Turing.counterLoopOrbit, Turing.counterLoopTM_run_done, Turing.counterLoopTM_run_escape, Turing.counterLoop_time_le, Turing.counterOverhead, Turing.counterOverhead_le_of_le, Turing.counterWord, Turing.tableBitsK}
\difficulty{6}
Let $k,n\in\mathbb N$, let $x$ be a finite binary input, let $M$ be a $k$-tape machine
with binary work tapes and states indexed by $\{0,\ldots,n\}$, and let $c$ be one of its
configurations whose state is present. Let $d$ be a configuration of the round machine
on $x$, at the live entry, representing $c$ and the transition table of $M$. Let $w$ be
a finite binary word, let $ct:\mathbb Z\to\{\text{absent},0,1\}$ be a tape and
$cp\in\mathbb Z$ a head position, and suppose that for every integer $p$ with $-1\le
p\le |w|$, the tape cell at $cp+p$ contains the symbol of $w$ at position $p$ (with
positions outside the word blank). If the binary numeral whose digits are the bits of
$w$ has value $t$, then there exist a time $T$, a configuration $e$, and a tape $ct'$
such that $T\le t\,50(L+1)+4t+2|w|+2$, where $L$ is the encoded transition-table length;
moreover, whenever the original machine has halted by a time $t'\le t$, one has $T\le
t'\,50(L+1)+4t'+2|w|$. The configuration $e$ represents the configuration of $M$ after
$t$ steps and has the same output as $d$. The tape $ct'$ agrees with $ct$ outside the
interval of cells from $cp$ through $cp+|w|-1$. At time $T$, the counter-loop machine,
started from the specified initial counter-loop configuration built from $d$, $ct$, and
$cp$, is at the configuration built from $e$ and $ct'$, with state determined by whether
$M$'s configuration after $t$ steps has a present state: it is the completed state if
present, and the escape state otherwise. At every time $u<T$, its state is neither the
completed state nor the escape state.
\end{theorem}

\begin{theorem}[Hosted simulation of a bounded loop]
\lean{Turing.StepBlock.hostedLoopTM_run}
\label{Turing.StepBlock.hostedLoopTM_run}
\leanok
\uses{Turing.Cfg, Turing.FinTM.VirtualTag, Turing.FinTM.bufferTape, Turing.MultiTapeTM, Turing.MultiTapeTM.runFrom, Turing.StepBlock.Ph, Turing.StepBlock.Rep, Turing.StepBlock.exit, Turing.StepBlock.hostedLoopTM, Turing.StepBlock.loopTM, Turing.StepBlock.loopTM_run, Turing.counterLoopCfg, Turing.tableBitsK, Turing.vhostCfg, Turing.vhostEmitTM, Turing.vhostEmitTM_runFrom}
\difficulty{4}
Let $M$ be a Turing machine with $k$ work tapes and state set $\mathrm{Fin}(n+1)$, and
let $c$ be an input configuration with a live state. Let $d$ be a configuration of the
compiled machine at its live entry that represents $c$. Let $w$ be a Boolean word, let
$ct:\mathbb{Z}\to\mathrm{Option}(\mathrm{Bool})$ be a tape, and let $cp\in\mathbb{Z}$.
Suppose $ct(cp+p)$ agrees with the buffer for $w$ for every integer $p$ with
$-1\le p\le |w|$, and that the binary value of $w$ is $t$. Then there are a time $T$, a
configuration $e$, and a tape $ct\prime$ such that
$T\le t(50(|\mathrm{tableBits}(M)|+1))+4t+2|w|+2$. For every $t\prime\le t$, if the
simulated machine has halted by time $t\prime$, then
$T\le t\prime(50(|\mathrm{tableBits}(M)|+1))+4t\prime+2|w|$. The configuration $e$
represents the simulated machine after $t$ steps and has the same output as $d$; $ct\prime$
agrees with $ct$ outside the interval $[cp,cp+|w|)$. At time $T$, the hosted machine
reaches the counter-loop configuration representing $e$ with tape $ct\prime$ and offset
$cp$; its state is the completed-loop exit if the simulated machine is still live, and the
escape state if it has halted. At every earlier time $u<T$, the hosted machine is in
neither the completed-loop exit state nor the escape state.
\end{theorem}

\begin{theorem}[Reading a hosted configuration]
\lean{Turing.StepBlock.hostedLoop_read}
\label{Turing.StepBlock.hostedLoop_read}
\leanok
\uses{Turing.Action, Turing.Cfg, Turing.Cfg.workTapeSymbols, Turing.CounterLoopState, Turing.StepBlock.Ph, Turing.StepBlock.Rep, Turing.StepBlock.outputSummary, Turing.StepBlock.summaryCells, Turing.StepBlock.summaryTape, Turing.counterLoopCfg, Turing.counter_cfg_eq, Turing.vhostBank, Turing.vhostCfg}
\difficulty{4}
Let $k,n$ be natural numbers, let $tr$ be a transition function for a $k$-work-tape
binary machine with states indexed by $\text{Fin}(n+1)$, and let $c'$ and $e$ be
configurations related by the compiled representation. For any loop-control state $q$,
tape content $ct:\mathbb{Z}\to\text{Option}(\text{Bool})$, offset $cp$, tag $b$, input-head
position $p$, and output prefix, the virtual-host configuration built from the loop
configuration at state $q$ has state $(q,b)$, input-head position $p$, output equal to
the prefix followed by $e$'s output, and the symbol read on the bank tape corresponding
to the summary tape is the output summary of $c'$'s output.
\end{theorem}
